\subsection{Définition et principes des SPA}

Une \textbf{SPA} (\textit{Single Page Application}) est une application web dont
l'interface est construite dynamiquement en JavaScript. Contrairement aux sites
classiques où chaque clic déclenche un rechargement HTML complet, la SPA charge le
shell de l'application une fois, puis met à jour le DOM (\textit{Document Object Model})
au fil des interactions utilisateur.

\medskip

Ce modèle repose sur le \textbf{routing côté client} : la bibliothèque React Router
associe chaque URL à un composant React (\texttt{/projects}, \texttt{/clients},
\texttt{/invoices}), sans solliciter le serveur pour le HTML. Seules les données
nécessaires sont récupérées via l'API REST.

\subsection{React : bibliothèque declarative}

\textbf{React}~\cite{react}, développé par Meta, est une bibliothèque JavaScript pour
construire des interfaces utilisateur par \textbf{composants} réutilisables. Chaque
composant encapsule structure (JSX), état et comportement.

\medskip

Les principes fondamentaux de React sont :

\begin{itemize}
  \item \textbf{Composants fonctionnels} : fonctions JavaScript retournant du JSX,
  composées hiérarchiquement (Sidebar, Dashboard, ProjectDetailPage, etc.).
  \item \textbf{État local} (\texttt{useState}) : données modifiables au sein d'un
  composant (formulaires, filtres, modales).
  \item \textbf{Effets de bord} (\texttt{useEffect}) : chargement de données API au
  montage du composant ou lors d'un changement de paramètre.
  \item \textbf{Context API} : partage d'état global (authentification, thème, langue)
  sans propager les props à travers toute l'arborescence.
\end{itemize}

Cette approche declarative simplifie la construction d'interfaces complexes : l'état
de l'interface est une fonction de l'état applicatif.

\subsection{Écosystème React du projet}

L'application CIVCO BTP s'appuie sur plusieurs bibliothèques complémentaires :

\begin{table}[H]
\caption{Composants de l'écosystème frontend}
\label{tab:react_ecosystem}
\centering
\small
\begin{tabular}{|l|l|p{7.5cm}|}
\hline
\textbf{Outil} & \textbf{Rôle} & \textbf{Usage dans le projet} \\
\hline
Vite~\cite{vite} & Build tool & Compilation rapide, HMR en développement \\
\hline
React Router & Routing & Navigation entre modules ERP \\
\hline
Axios & Client HTTP & Appels API REST avec intercepteurs \\
\hline
Tailwind CSS~\cite{tailwindcss} & Styles utilitaires & Interface sombre cohérente \\
\hline
Recharts & Graphiques & Tableau de bord, KPIs \\
\hline
Leaflet & Cartographie & Carte des projets / chantiers \\
\hline
\end{tabular}
\end{table}

\subsection{Communication avec le backend}

Le \textit{frontend} React consomme l'API Laravel via des modules dédiés
(\texttt{api/projects.js}, \texttt{api/clients.js}, etc.). Un client Axios centralisé
configure :

\begin{itemize}
  \item l'envoi automatique des cookies de session (Sanctum) ;
  \item le jeton CSRF pour les requêtes \texttt{POST}, \texttt{PUT} et \texttt{DELETE} ;
  \item la gestion des erreurs \texttt{401} (redirection vers la page de connexion).
\end{itemize}

Ce patron \textit{API-first} garantit que toute action utilisateur (créer un devis,
archiver un client, assigner une tâche) transite par la couche métier Laravel, où les
règles de validation et de permissions sont appliquées.

\begin{figure}[H]
\centering
\resizebox{0.92\textwidth}{!}{%
\begin{tikzpicture}[
  box/.style={rectangle, rounded corners=4pt, fill=ensamBlue!12, draw=ensamBlue,
              minimum width=3.2cm, minimum height=1.0cm, align=center,
              font=\small\bfseries},
  arr/.style={-{Stealth[length=6pt]}, thick, draw=gray!70},
  lbl/.style={font=\footnotesize, text=gray!70}
]
  \node[box] (browser) at (0, 0)   {Navigateur\\React SPA};
  \node[box] (api)     at (5.5, 0) {API Laravel\\REST + Sanctum};
  \node[box] (db)      at (11, 0)  {Base de\\données};

  \draw[arr] (browser) -- (api) node[midway, above, lbl] {JSON / HTTPS};
  \draw[arr] (api)     -- (db)  node[midway, above, lbl] {Eloquent / SQL};
\end{tikzpicture}
}%
\caption{Flux de communication React $\rightarrow$ Laravel $\rightarrow$ SGBD}
\label{fig:react_flow}
\end{figure}

\subsection{Comparaison React / Vue / Angular}

\begin{table}[H]
\caption{Comparaison des principaux frameworks frontend (2025-2026)}
\label{tab:spa_compare}
\centering
\small
\begin{tabular}{|l|p{3.5cm}|p{3.5cm}|p{3.5cm}|}
\hline
\textbf{Critère} & \textbf{React} & \textbf{Vue.js} & \textbf{Angular} \\
\hline
Flexibilité & Élevée (bibliothèque) & Élevée (progressif) & Opinionated (framework complet) \\
\hline
Courbe d'apprentissage & Modérée & Faible & Élevée \\
\hline
Marché de l'emploi & Très large & Large & Large (entreprise) \\
\hline
Écosystème SPA & React Router, Vite & Vue Router, Vite & Router intégré \\
\hline
\end{tabular}
\end{table}

React est retenu pour sa flexibilité, son adoption industrielle et la disponibilité de
bibliothèques adaptées aux tableaux de bord métier (graphiques, cartes, composants
modulaires), comme le montre le Tableau~\ref{tab:spa_compare}.

\subsection{Organisation du code frontend}

Le code React est structuré par \textbf{features} (\texttt{features/projects/},
\texttt{features/clients/}, \texttt{features/invoices/}), chacune regroupant pages,
composants et utilitaires propres au module. Les composants transverses (Sidebar,
PermissionGate, modales) vivent dans \texttt{components/}. Cette organisation par
domaine facilite la navigation dans le code et l'ajout de nouvelles fonctionnalités
sans impacter les modules existants.

\subsection{Interface utilisateur et Tailwind CSS}

L'interface adopte un design sombre professionnel, cohérent sur tous les modules.
\textbf{Tailwind CSS}~\cite{tailwindcss} est un framework CSS utilitaire : les classes
(\texttt{flex}, \texttt{rounded-lg}, \texttt{bg-gray-800}) s'appliquent directement
dans le JSX. Cette approche accélère le prototypage et garantit une cohérence visuelle
sur les écrans de l'application (listes, fiches détail, modales, tableaux de bord).

\medskip

Le \textbf{tableau de bord} illustre l'intérêt de React pour la visualisation métier :
graphiques (Recharts), widgets configurables, indicateurs \textbf{KPI} et accès rapide
aux projets en cours. La carte des chantiers (Leaflet) offre une vue géographique des
projets, pertinente pour une entreprise disposant de plusieurs sites simultanés.

\subsection{Internationalisation}

L'application supporte le \textbf{français} (langue par défaut) et l'\textbf{anglais}
via un contexte de langue et des fichiers de traduction (\texttt{locales/fr.js},
\texttt{locales/en.js}). Cette préparation facilite l'usage de l'outil par des
collaborateurs anglophones ou une future extension à d'autres contextes linguistiques.
